\bdchap{ch:intro}{引论}
\SourcePage{3}

\bdp{0.1} 设 $X$ 是 $\mathbb{C}$ 上一条亏格 $g>1$ 的连通光滑射影曲线，$G$ 是 $\mathbb{C}$ 上的半单群，$\mathfrak{g}$ 是 $G$ 的李代数。记 $\operatorname{Bun}_G=\operatorname{Bun}_G(X)$ 为 $X$ 上 $G$-丛的模栈。Hitchin 在 [Hit87] 中定义了 $\operatorname{Bun}_G(X)$ 的余切栈上一个值得注意的代数 $\mathfrak{z}^{cl}=\mathfrak{z}^{cl}(X)$，其元素是相互泊松交换的函数。

\bdp{0.2} 本注记证明以下结论：

\begin{enumerate}[label=(\alph*),leftmargin=2.2em]
\item Hitchin 构造具有自然的量子化。确切地说，我们定义 $\operatorname{Bun}_G$ 上一族扭微分算子构成的交换环 $\mathfrak{z}=\mathfrak{z}(X)$，使得 $\mathfrak{z}$ 中算子的象征恰好构成 Hitchin 哈密顿量的环 $\mathfrak{z}^{cl}$。这里「扭」是指所考虑的微分算子作用在典范丛 $\omega_{\operatorname{Bun}_G}$ 的一个平方根上。扭是本质的：众所周知，$\operatorname{Bun}_G$ 上唯一的整体非扭微分算子是乘以局部常值函数。

\item $\mathfrak{z}$ 的谱与 $^{L}\mathfrak{g}$-oper 的模空间有典范等同，后者是 $X$ 上不可约 $(^{L}G)_{ad}$-局部系统的模空间的一个（Lagrange）子空间。这里 $^{L}G$ 是 $G$ 的朗兰兹对偶，$^{L}\mathfrak{g}$ 是其李代数，$(^{L}G)_{ad}$ 是伴随群；关于 oper 的简短说明见 0.3。

\item 对 $^{L}\mathfrak{g}$-oper $\mathfrak{F}$，记 $N_{\mathfrak{F}}$ 为扭微分算子层商去由极大理想 $\mathfrak{m}_{\mathfrak{F}}\subset\mathfrak{z}$ 生成的左理想所得的商。它是 $\operatorname{Bun}_G$ 上一个非零的完整扭 $\mathcal{D}$-模。

\item 对 $^{L}G$-oper $\mathfrak{F}$，可以配置 $\operatorname{Bun}_G$ 上一个通常的（非扭）$\mathcal{D}$-模 $M_{\mathfrak{F}}$。若 $G$ 单连通，则 $M_{\mathfrak{F}}$ 同构于 $\omega_{\operatorname{Bun}_G}^{-1/2}\otimes N_{\mathfrak{F}}$（单连通时 $\omega_{\operatorname{Bun}_G}^{1/2}$ 唯一；另一方面 $N_{\mathfrak{F}}$ 有意义，因为此时 $^{L}G$-oper 与 $^{L}\mathfrak{g}$-oper 没有区别）。一般地，$M_{\mathfrak{F}}:=\lambda_{\mathfrak{F}}^{-1}\otimes N_{\bar F}$，其中 $\bar F$ 是对应于 $\mathfrak{F}$ 的 $^{L}\mathfrak{g}$-oper，$\lambda_{\mathfrak{F}}$ 是 $\operatorname{Bun}_G$ 上某个带扭 $\mathcal{D}$-模结构的可逆层（见 5.1.1）。$\lambda_{\mathfrak{F}}$ 的同构类只依赖于 $\mathfrak{F}$ 在 $^{L}G$-oper 模空间中的连通分支。

\item 主定理：$M_{\mathfrak{F}}$ 是以 $\mathfrak{F}$ 为特征值的 Hecke 特征层（确切的陈述见 ???）。
% [Tutor Note: 原书 p.4 此处的交叉引用印作「???」，是 LaTeX 未解析的引用；底本如此，保留原样。]
换言之，在朗兰兹的意义下 $M_{\mathfrak{F}}$ 对应于局部系统 $\mathfrak{F}$。
\end{enumerate}

\SourcePage{4}

\bdp{0.3} oper 这个概念（而非名称）是相当为人熟知的（例如，相应的局部对象已在 [DS85] 中研究过）。光滑曲线 $Y$ 上的 $G$-oper 是一个 $G$-局部系统（即带联络的 $G$-丛），并带有某种附加结构（见 3.1.3）。若 $G=\mathrm{SL}_n$（此时我们处理的是向量空间的局部系统），则 oper 结构是一个子丛的完全旗，它在 $Y$ 的每一点都满足 Griffiths 横截性条件与相应的非退化条件。$\mathrm{PSL}_2$-oper 等同于 $Y$ 上的射影联络，即 $Y$ 上的 Sturm-Liouville 算子（见 [Del70]）。按定义，$\mathfrak{g}$-oper 是伴随群 $G_{ad}$ 的 oper。若 $Y$ 完备且亏格为正，则一个局部系统至多承载一个 oper 结构，因此可以把 oper 视为特殊的局部系统。

\bdp{0.4} 0.2 中的整体构造与陈述有局部对应物，后者起主要作用。(a)、(b) 的局部版本是下述典范同构：$\mathfrak{g}((t))$ 的临界扭（完备化）包络代数的中心的谱，与去心圆盘 $\operatorname{Spec}\mathbb{C}((t))$ 上 $^{L}\mathfrak{g}$-oper 的模空间之间的典范同构。该同构由 Feigin 与 Frenkel [FF92] 建立，它是 $\mathfrak{g}$ 与 $^{L}\mathfrak{g}$ 的 W-代数之间一个值得注意的对称性的特化。我们不知道这一「双重量子化」的图景能否整体化。0.2(c)、(d) 的局部版本本质上给出 Feigin--Frenkel 同构的另一种构造，它基于 Bruhat--Tits 仿射 Grassmann 流形的几何。这里关键是一个关于某些临界扭 $\mathcal{D}$-模的上同调的消灭定理（对「次临界」扭的平行结果已在 [KT95] 中证明）。

\bdp{0.5} 本注记只包含主要结果证明的概要。若干技术性结果只作陈述而不给证明。详细阐述将在后续论文中给出。

\bdp{0.6} 我们想指出，E. Witten 独立发现了 0.2(a--d) 的想法，并猜测了 0.2(e)。据我们所知，他没有就此发表任何东西。

\bdp{0.7} 本文结果的一个较弱版本已在 [BD96] 中宣布。

\SourcePage{5}

\bdp{0.8} 作者感谢 P. Deligne、V. Ginzburg、B. Feigin 与 E. Frenkel 的有益讨论。我们也感谢普林斯顿高等研究院（Institute for Advanced Study）的接待。我们衷心感谢 R. Becker、W. Snow、D. Phares 与 S. Fryntova 仔细录入手稿。
